% Generated draft: requires mathematical and source review.
\documentclass[11pt]{article}
\usepackage[margin=25mm]{geometry}
\usepackage{fontspec}
\setmainfont{Latin Modern Roman}
\usepackage{amsmath,amssymb,mathtools,unicode-math}
\setmathfont{Latin Modern Math}
\usepackage{enumitem,needspace,xurl,hyperref}
\setlength{\parindent}{0pt}
\setlength{\parskip}{5pt}
\setlength{\emergencystretch}{3em}
\newcommand{\N}{\mathbb{N}}
\newcommand{\Z}{\mathbb{Z}}
\newcommand{\Q}{\mathbb{Q}}
\newcommand{\R}{\mathbb{R}}
\newcommand{\C}{\mathbb{C}}
\newcommand{\F}{\mathbb{F}}
\newcommand{\A}{\mathbb{A}}
\newcommand{\PP}{\mathbb{P}}
\newcommand{\E}{\mathbb{E}}
\newcommand{\eps}{\varepsilon}
\newcommand{\dd}{\,\mathrm{d}}
\begin{document}
\section*{UnsolvedMath ID: 700832}
\textbf{Source problem code:} catalogue entry
\subsection{Definitions and mathematical statement}
A lattice-ordered group is a group \(G\) equipped with a partial order \(\leq\) such that every pair of elements has a meet \(\wedge\) and a join \(\vee\), and
\[
x\leq y\ \Longrightarrow\ axb\leq ayb
\qquad(a,b,x,y\in G).
\]
Define the lower central series of the underlying group by \(\gamma_1(G)=G\) and
\[
\gamma_{n+1}(G)=\langle [u,v]:u\in\gamma_n(G),\ v\in G\rangle,
\qquad [u,v]=u^{-1}v^{-1}uv.
\]
The group is nilpotent if \(\gamma_{c+1}(G)=\{1\}\) for some integer \(c\geq 0\).

For a finite tuple \(S=(s_1,\ldots,s_r)\), generation as a lattice-ordered group means that every element is the value of a finite term in \(S\) and \(1\), using multiplication, inversion, meet, and join. Generation as an abstract group means that multiplication and inversion alone suffice. The supplied statement does not explicitly distinguish these meanings of finite generation.

For a fixed lattice-ordered group \(G\) generated by \(S\) as a lattice-ordered group, define its conjugacy problem on lattice-group terms as follows: given two finite terms \(t,u\) in \(S\), decide whether
\[
\exists g\in G\quad g^{-1}t(S)g=u(S).
\]
Solubility means that an algorithm halts on every such pair and gives the correct answer.

\textbf{Question attributed to A. M. W. Glass (1992).} Is the conjugacy problem soluble for every nilpotent finitely generated lattice-ordered group?

Under the interpretation that finite generation is in the lattice-ordered group signature and inputs are lattice-group terms, this asks whether, for every such \(G\) and finite generating tuple \(S\), there exists an algorithm deciding the displayed condition for all \(t,u\). This is an explicit interpretation, not a verified clarification of the source. If finite generation instead refers to the underlying abstract group, the class of groups changes and inputs may be taken to be ordinary group words. The record also does not specify an effective presentation or whether a single algorithm must work uniformly from presentations of the groups. These qualifications remain unresolved.
\subsection{Short English statement}
Glass asks whether conjugacy can always be decided in a nilpotent finitely generated lattice-ordered group: given two represented elements, determine whether one is a conjugate of the other. The supplied record leaves the meaning of finite generation, the input representation, and uniformity requirements unspecified. Its \emph{open} label is supplied registry metadata, not an independently verified assessment.
\subsection{Sources}
\begin{itemize}
\item Local UnsolvedMath registry, permanent ID \(700832\), exact source problem code \texttt{catalogue entry}, title \emph{Kourovka Notebook Problem 12.12}. The supplied record labels the problem \emph{open} and attributes it to A. M. W. Glass, proposed in 1992.
\item Supplied source attribution: \emph{The Kourovka Notebook (Unsolved Problems in Group Theory)}, 12th issue (1992), original problem number 12.12; also identified in the record as the 21st issue edition (2026), \texttt{arXiv:1401.0300v48}. The supplied collection description names E. I. Khukhro and V. D. Mazurov as editors.
\item These bibliographic details and the problem statement are supplied source data. The cited edition and its contents have not been independently consulted or verified for this draft.
\end{itemize}
\end{document}
